\section{Subgroup labels and simultaneous symmetry}
\label{sec:grouped_label_symmetric_support}

For a permutation action $\phi$ of $S_r$, write $\Pi_\phi$ for the
projection onto its fixed subspace and
$F_\phi=\sum_\lambda\log(\dim\lambda)\pi_\phi^\lambda$.

\begin{theorem}[Complementary subgroup labels]
\label{thm:complementary_subgroup_label_entropies}
\lean{PermutationRepresentation.labelEntropy_mul_symProj_of_eq_mul}
Let $\phi,\chi,\psi$ be actions of $S_r$ on the same finite
coordinate set. If $\psi(\sigma)=\phi(\sigma)\chi(\sigma)$ for
every $\sigma$, then $F_\phi\Pi_\psi=F_\chi\Pi_\psi$.
This is the complementary-label argument of
\cite[Lemma~6.1(2)]{OpenAI2026AreaLaw}.
\end{theorem}
\begin{proof}
On a vector fixed by $\psi$, the action of $\chi(\sigma)$ equals
that of $\phi(\sigma^{-1})$. The coefficients of every central
idempotent of $\C[S_r]$ are invariant under inversion. Apply this
identity to each central idempotent, then sum with coefficients
$\log\dim\lambda$.
\end{proof}

Let $\mathcal H=\bigotimes_{v\in V}\C^{I_v}$ with $V$ and all
$I_v$ finite. For a homomorphism $\theta:S_r\to S_k$, write
$F_S^\theta$ for the label entropy of the induced copy action on
$S\subseteq V$, and $\Pi_T^\theta$ for simultaneous symmetry on
$T$ under that subgroup.

\begin{theorem}[A signed subgroup entropy on its symmetric range]
\label{thm:signed_subgroup_label_symmetric_nonnegative}
\lean{TensorPower.commute_subgroup_labelObservable_symProj_of_subset,
  TensorPower.subgroup_signedLabelEntropy_symProj_nonneg}
\uses{thm:complementary_subgroup_label_entropies}
Every real central label observable on $S$ commutes with
$\Pi_T^\theta$ when $S\subseteq T$. If $P,Y,F$ are pairwise
disjoint and cover $V$, then, with
$G=F_P^\theta+F_F^\theta-F_Y^\theta$ and
$\Pi=\Pi_V^\theta$,
\begin{align}
 G\Pi&=\Pi G,&\Pi G\Pi&\ge0.
\end{align}
The signed operator need not be nonnegative on the entire space.
\end{theorem}
\begin{proof}
Centrality on a smaller region and commutation with its disjoint
complement give the first assertion. Complementary labels replace
$F_Y^\theta\Pi$ by $F_{P\cup F}^\theta\Pi$. Thus
$G\Pi=D\Pi$, where
$D=F_P^\theta+F_F^\theta-F_{P\cup F}^\theta\ge0$ is the actual
merge deficit. Hence $\Pi G\Pi=\Pi D\Pi\ge0$.
\end{proof}

\begin{theorem}[Commutation between disjoint groups of copies]
\label{thm:grouped_copy_label_commutation}
\lean{TensorPower.commute_subsystemPerm_of_commute,
  TensorPower.commute_grouped_labelObservables,
  TensorPower.commute_good_labelObservable_bad_symProj}
Commuting permutations of the copies induce commuting actions on
any two physical subsystems, including overlapping subsystems.
Consequently, after a division of the copies into groups of sizes
$m$ and $r$, every real label observable of the first group commutes
with every such observable of the second group. It also commutes
with the simultaneous symmetric projection of the second group.
These are the subgroup commutations used in
\cite[Section~8, lines 454--480]{OpenAI2026AreaLaw}.
\end{theorem}
\begin{proof}
On the intersection of the two subsystems, the induced actions
commute because the underlying copy permutations commute. Outside
the intersection, at most one of the actions moves a coordinate.
Pass from these permutation identities to their group-algebra
images and to the average defining the symmetric projection.
\end{proof}
